\section{EigenLayer Restaking Background}
\label{sec:background}

\subsection{From Staking to Restaking}

Ethereum validators bond 32~ETH to participate in consensus.  This
stake secures one service\,---\,Ethereum's own block production and
finality\,---\,and is slashable for protocol-level
faults~\cite{buterin2014ethereum}.  EigenLayer's central observation
is that this same capital can simultaneously back \emph{additional}
services if the staker explicitly opts in to additional slashing
conditions for them~\cite{eigenlayer2023}.  The EigenLayer contracts
expose a \texttt{Strategy} primitive that lets a staker delegate
their bonded ETH (or ETH-derivative LSTs) to one or more Operators
who run AVS clients.  Slashing on any registered AVS reduces the
underlying restaked balance, and the shortfall is socialized to the
delegators.

\subsection{The Four AVS Roles}

A typical AVS implemented on top of EigenLayer comprises four roles.
The \emph{ServiceManager} is a Solidity contract that the AVS
operator deploys to register itself with the EigenLayer core, declare
its slashing rules, and gate operator opt-in.  The
\emph{TaskManager} is the AVS-specific contract that emits tasks and
verifies responses; in our system it owns the ETH/USD price-quote
schedule.  \emph{Operators} are off-chain agents (Go binaries in our
case) that listen for new tasks, perform the AVS-defined computation,
sign the response with their BLS private key, and forward the
signature to an \emph{Aggregator}.  The Aggregator collects signatures
until the quorum threshold is met, combines them into a single
BLS-aggregated signature, and submits the final response on chain
where the TaskManager verifies it~\cite{layrlabs_isavs,eigenlabs_avsbook}.

\subsection{Alternatives and Shared Security}

EigenLayer is not the only shared-security
design.  Symbiotic~\cite{symbiotic2024} offers a
permissionless variant where any ERC-20 collateral can secure any
network, with slashing logic externalized to the network's own
contracts.  Karak~\cite{karak2024} provides a similar restaking layer
across multiple L1s with a focus on developer ergonomics.  All three
designs share the same fundamental risk:
\emph{cascading slashing}\,---\,a single bug in any participating
AVS can cause coordinated losses across the entire restaker
set.  Durvasula and Roughgarden formalize this risk and derive
sufficient conditions on individual-AVS slashing magnitudes that
guarantee whole-network solvency~\cite{durvasula2024robust}.  Our
work targets EigenLayer because of its mature tooling and reference
template, but the design lessons transfer to the other two.

\begin{figure}[t]
  \centering
  % \includegraphics[width=0.95\columnwidth]{figures/fig_eigenlayer_overview.pdf}
  \framebox[0.95\columnwidth][c]{\raisebox{1.2cm}{\textit{[placeholder: EigenLayer architecture diagram --- M4 to add]}}}
  \caption{High-level data flow between the AVS contracts, the
           registered Operators (whose ETH is restaked), and the
           off-chain Aggregator and Challenger services.  Adapted
           from~\cite{eigenlabs_avsbook}.}
  \label{fig:eigenlayer-overview}
\end{figure}
